\documentclass[12pt,a4paper]{article}

\usepackage[T2A]{fontenc}
\usepackage[utf8]{inputenc}
\usepackage[russian]{babel}
\usepackage{amsmath,amssymb,amsthm,mathtools}
\usepackage{enumitem}
\usepackage{listings}
\usepackage{xcolor}
\usepackage{tikz}
\usepackage{array}
\usepackage{booktabs}
\usepackage{hyperref}

\DeclareMathOperator{\lfp}{lfp}
\DeclareMathOperator{\gfp}{gfp}

\newcommand{\BigSqcap}{\mathop{\mathchoice
  {\vcenter{\hbox{\Large$\sqcap$}}}
  {\vcenter{\hbox{\large$\sqcap$}}}
  {\vcenter{\hbox{\normalsize$\sqcap$}}}
  {\vcenter{\hbox{\small$\sqcap$}}}}\displaylimits}

\newcommand{\BigSqcup}{\mathop{\mathchoice
  {\vcenter{\hbox{\Large$\sqcup$}}}
  {\vcenter{\hbox{\large$\sqcup$}}}
  {\vcenter{\hbox{\normalsize$\sqcup$}}}
  {\vcenter{\hbox{\small$\sqcup$}}}}\displaylimits}

\usetikzlibrary{arrows.meta,positioning,shapes.geometric}

\newtheorem{definition}{Определение}
\newtheorem{theorem}{Теорема}
\newtheorem{lemma}{Лемма}
\newtheorem{remark}{Замечание}
\newtheorem{example}{Пример}

\lstdefinestyle{code}{
  basicstyle=\ttfamily\small,
  keywordstyle=\color{blue!70!black},
  commentstyle=\color{green!40!black},
  stringstyle=\color{red!60!black},
  numbers=left,
  numberstyle=\tiny,
  stepnumber=1,
  numbersep=8pt,
  showstringspaces=false,
  tabsize=2,
  breaklines=true,
  frame=single,
  inputencoding=utf8,
  extendedchars=true,
  literate=
    {а}{{\selectfont\char224}}1 {б}{{\selectfont\char225}}1 {в}{{\selectfont\char226}}1
    {г}{{\selectfont\char227}}1 {д}{{\selectfont\char228}}1 {е}{{\selectfont\char229}}1
    {ё}{{\"e}}1 {ж}{{\selectfont\char230}}1 {з}{{\selectfont\char231}}1
    {и}{{\selectfont\char232}}1 {й}{{\selectfont\char233}}1 {к}{{\selectfont\char234}}1
    {л}{{\selectfont\char235}}1 {м}{{\selectfont\char236}}1 {н}{{\selectfont\char237}}1
    {о}{{\selectfont\char238}}1 {п}{{\selectfont\char239}}1 {р}{{\selectfont\char240}}1
    {с}{{\selectfont\char241}}1 {т}{{\selectfont\char242}}1 {у}{{\selectfont\char243}}1
    {ф}{{\selectfont\char244}}1 {х}{{\selectfont\char245}}1 {ц}{{\selectfont\char246}}1
    {ч}{{\selectfont\char247}}1 {ш}{{\selectfont\char248}}1 {щ}{{\selectfont\char249}}1
    {ъ}{{\selectfont\char250}}1 {ы}{{\selectfont\char251}}1 {ь}{{\selectfont\char252}}1
    {э}{{\selectfont\char253}}1 {ю}{{\selectfont\char254}}1 {я}{{\selectfont\char255}}1
    {А}{{\selectfont\char192}}1 {Б}{{\selectfont\char193}}1 {В}{{\selectfont\char194}}1
    {Г}{{\selectfont\char195}}1 {Д}{{\selectfont\char196}}1 {Е}{{\selectfont\char197}}1
    {Ё}{{\"E}}1 {Ж}{{\selectfont\char198}}1 {З}{{\selectfont\char199}}1
    {И}{{\selectfont\char200}}1 {Й}{{\selectfont\char201}}1 {К}{{\selectfont\char202}}1
    {Л}{{\selectfont\char203}}1 {М}{{\selectfont\char204}}1 {Н}{{\selectfont\char205}}1
    {О}{{\selectfont\char206}}1 {П}{{\selectfont\char207}}1 {Р}{{\selectfont\char208}}1
    {С}{{\selectfont\char209}}1 {Т}{{\selectfont\char210}}1 {У}{{\selectfont\char211}}1
    {Ф}{{\selectfont\char212}}1 {Х}{{\selectfont\char213}}1 {Ц}{{\selectfont\char214}}1
    {Ч}{{\selectfont\char215}}1 {Ш}{{\selectfont\char216}}1 {Щ}{{\selectfont\char217}}1
    {Ъ}{{\selectfont\char218}}1 {Ы}{{\selectfont\char219}}1 {Ь}{{\selectfont\char220}}1
    {Э}{{\selectfont\char221}}1 {Ю}{{\selectfont\char222}}1 {Я}{{\selectfont\char223}}1
}

\lstset{style=code}

\title{Билет №41 \\ \small Принципы статической и динамической проверки программного кода. (СпБПУ)}
\author{Сериков Владислав Александрович}
\date{13 мая 2026}

\begin{document}

\maketitle
\tableofcontents
\newpage

\section{Общая постановка задачи проверки программ}

Проверка программного кода направлена на обнаружение дефектов, уязвимостей и несоответствий программы заданным требованиям. В широком смысле различают два фундаментальных подхода:

\begin{itemize}
  \item \textbf{статическая проверка} --- анализ программы без её выполнения;
  \item \textbf{динамическая проверка} --- анализ поведения программы во время выполнения.
\end{itemize}

Оба подхода дополняют друг друга. Статическая проверка может исследовать множество возможных путей выполнения, включая те, которые редко возникают на практике, но обычно страдает от ложных срабатываний. Динамическая проверка, наоборот, наблюдает реальные выполнения программы и поэтому часто даёт более конкретные свидетельства дефектов, но не может доказать отсутствие ошибок во всех возможных выполнениях, если не покрыты все состояния и входы.

\begin{definition}
\textbf{Дефектом} программы называется ошибка в коде, спецификации, архитектуре или предположениях разработчика, способная привести к неправильному поведению программы.
\end{definition}

\begin{definition}
\textbf{Отказом} называется наблюдаемое нарушение требуемого поведения программы во время её выполнения.
\end{definition}

\begin{definition}
\textbf{Проверкой} программы называется процесс получения свидетельств о корректности или некорректности программы относительно некоторого набора требований, спецификаций или свойств безопасности.
\end{definition}

Типичные свойства, проверяемые статически и динамически:

\begin{itemize}
  \item отсутствие обращения по нулевому указателю;
  \item отсутствие выхода за границы массива;
  \item отсутствие использования неинициализированных переменных;
  \item отсутствие утечек памяти и ресурсов;
  \item отсутствие гонок данных;
  \item соответствие интерфейсным контрактам;
  \item достижимость или недостижимость определённых участков кода;
  \item соблюдение стандартов кодирования;
  \item выполнение функциональных требований.
\end{itemize}

\section{Статическая проверка программного кода}

\subsection{Сущность статической проверки}

\begin{definition}
\textbf{Статическая проверка программного кода} --- это анализ текста программы, её промежуточного представления, байткода, объектного кода или модели без запуска программы на конкретных входных данных.
\end{definition}

Статическая проверка может выполняться на разных уровнях представления программы:

\begin{itemize}
  \item исходный текст;
  \item абстрактное синтаксическое дерево;
  \item таблицы символов;
  \item граф потока управления;
  \item граф потока данных;
  \item промежуточное представление компилятора;
  \item байткод;
  \item машинный код.
\end{itemize}

Типовой конвейер статического анализа:

\begin{center}
\begin{tikzpicture}[
  node distance=1.4cm,
  block/.style={rectangle, draw, rounded corners, align=center, minimum width=3.2cm, minimum height=0.8cm},
  arrow/.style={-{Latex[length=3mm]}}
]
\node[block] (src) {Исходный код};
\node[block, right=of src] (lex) {Лексический\\анализ};
\node[block, right=of lex] (parse) {Синтаксический\\анализ};
\node[block, below=of parse] (sem) {Семантический\\анализ};
\node[block, left=of sem] (ir) {Промежуточное\\представление};
\node[block, left=of ir] (checks) {Проверки\\и отчёты};

\draw[arrow] (src) -- (lex);
\draw[arrow] (lex) -- (parse);
\draw[arrow] (parse) -- (sem);
\draw[arrow] (sem) -- (ir);
\draw[arrow] (ir) -- (checks);
\end{tikzpicture}
\end{center}

\subsection{Основные виды статической проверки}

\subsubsection{Лексическая и синтаксическая проверка}

Лексическая проверка разбивает текст программы на токены: идентификаторы, ключевые слова, литералы, операторы, разделители.

Синтаксическая проверка строит дерево разбора или абстрактное синтаксическое дерево и проверяет, принадлежит ли программа грамматике языка.

\begin{example}
В языке C выражение
\begin{lstlisting}[language=C]
int x = ;
\end{lstlisting}
синтаксически некорректно, поскольку после оператора присваивания отсутствует выражение.
\end{example}

\subsubsection{Семантическая проверка}

Семантическая проверка обнаруживает ошибки, которые не являются чисто синтаксическими:

\begin{itemize}
  \item использование необъявленной переменной;
  \item несовместимость типов;
  \item неверное число аргументов функции;
  \item нарушение правил областей видимости;
  \item неправильное использование модификаторов доступа;
  \item невозможность неявного преобразования типа.
\end{itemize}

\begin{example}
\begin{lstlisting}[language=C]
int x = "hello";
\end{lstlisting}
Синтаксически эта запись допустима, но семантически некорректна: строковый литерал не может быть непосредственно присвоен переменной типа \texttt{int}.
\end{example}

\subsubsection{Проверка типов}

\begin{definition}
\textbf{Система типов} --- это формальная система, приписывающая выражениям типы и ограничивающая допустимые операции над значениями.
\end{definition}

Типизация является одной из наиболее распространённых форм статической проверки.

\begin{example}
Если переменная \texttt{x} имеет тип \texttt{int}, а переменная \texttt{s} имеет тип \texttt{string}, то выражение
\[
x + s
\]
может быть запрещено системой типов, если язык не определяет сложение целого числа и строки.
\end{example}

\subsubsection{Анализ потока управления}

\begin{definition}
\textbf{Граф потока управления} программы --- это ориентированный граф
\[
CFG = (V, E),
\]
где вершины $V$ соответствуют базовым блокам или операторам программы, а дуги $E$ соответствуют возможным переходам управления.
\end{definition}

\begin{definition}
\textbf{Базовый блок} --- это максимальная последовательность команд, в которую управление входит только в начало и выходит только из конца.
\end{definition}

Пример:

\begin{lstlisting}[language=C]
int f(int x) {
    if (x > 0)
        return 1;
    else
        return -1;
}
\end{lstlisting}

Граф потока управления:

\begin{center}
\begin{tikzpicture}[
  node distance=1.2cm,
  block/.style={rectangle, draw, rounded corners, align=center, minimum width=3cm, minimum height=0.8cm},
  arrow/.style={-{Latex[length=3mm]}}
]
\node[block] (entry) {Вход};
\node[block, below=of entry] (cond) {$x > 0$?};
\node[block, below left=of cond] (ret1) {\texttt{return 1}};
\node[block, below right=of cond] (ret2) {\texttt{return -1}};
\node[block, below=2cm of cond] (exit) {Выход};

\draw[arrow] (entry) -- (cond);
\draw[arrow] (cond) -- node[left] {да} (ret1);
\draw[arrow] (cond) -- node[right] {нет} (ret2);
\draw[arrow] (ret1) -- (exit);
\draw[arrow] (ret2) -- (exit);
\end{tikzpicture}
\end{center}

На основе графа потока управления можно обнаруживать:

\begin{itemize}
  \item недостижимый код;
  \item бесконечные циклы;
  \item отсутствие возврата значения из функции;
  \item потенциально неисполненные ветви;
  \item сложность программы.
\end{itemize}

\subsubsection{Анализ потока данных}

\begin{definition}
\textbf{Анализ потока данных} --- это семейство методов статического анализа, вычисляющих информацию о возможных значениях и состояниях программы в точках графа потока управления.
\end{definition}

Классические задачи анализа потока данных:

\begin{itemize}
  \item достижимые определения;
  \item живые переменные;
  \item доступные выражения;
  \item константное распространение;
  \item анализ инициализированности;
  \item анализ указателей и алиасов.
\end{itemize}

\subsection{Анализ достижимых определений}

\begin{definition}
Определение переменной $x$ в точке $p$ программы \textbf{достигает} точки $q$, если существует путь от $p$ к $q$, на котором переменная $x$ не переопределяется.
\end{definition}

Для каждого базового блока $B$ вводятся множества:

\begin{itemize}
  \item $GEN[B]$ --- определения, порождённые блоком $B$;
  \item $KILL[B]$ --- определения, уничтоженные блоком $B$;
  \item $IN[B]$ --- определения, достигающие входа в блок $B$;
  \item $OUT[B]$ --- определения, достигающие выхода из блока $B$.
\end{itemize}

Уравнения анализа:

\[
IN[B] = \bigcup_{P \in pred(B)} OUT[P],
\]

\[
OUT[B] = GEN[B] \cup (IN[B] \setminus KILL[B]).
\]

\subsubsection{Алгоритм итерационного анализа достижимых определений}

\textbf{Вход:} граф потока управления $CFG$, множества $GEN[B]$ и $KILL[B]$ для каждого блока.

\textbf{Выход:} множества $IN[B]$ и $OUT[B]$.

\begin{enumerate}
  \item Для всех блоков $B$ положить:
  \[
  IN[B] := \varnothing,\qquad OUT[B] := \varnothing.
  \]
  \item Повторять, пока хотя бы одно множество изменяется:
  \begin{enumerate}
    \item для каждого блока $B$ вычислить
    \[
    IN[B] := \bigcup_{P \in pred(B)} OUT[P];
    \]
    \item вычислить
    \[
    OUT[B] := GEN[B] \cup (IN[B] \setminus KILL[B]).
    \]
  \end{enumerate}
  \item Вернуть найденные множества.
\end{enumerate}

\subsubsection{Минимальный пример}

Рассмотрим программу:

\begin{lstlisting}[language=C]
1: x = 1;
2: if (c)
3:     x = 2;
4: y = x;
\end{lstlisting}

Определения:

\[
d_1: x = 1,\qquad d_2: x = 2.
\]

Блоки:

\[
B_1: x=1,\quad B_2: if(c),\quad B_3: x=2,\quad B_4: y=x.
\]

Граф:

\begin{center}
\begin{tikzpicture}[
  node distance=1.1cm,
  block/.style={rectangle, draw, rounded corners, align=center, minimum width=2.8cm, minimum height=0.8cm},
  arrow/.style={-{Latex[length=3mm]}}
]
\node[block] (b1) {$B_1: x=1$};
\node[block, below=of b1] (b2) {$B_2: if(c)$};
\node[block, below left=of b2] (b3) {$B_3: x=2$};
\node[block, below right=of b2] (b4) {$B_4: y=x$};

\draw[arrow] (b1) -- (b2);
\draw[arrow] (b2) -- node[left] {true} (b3);
\draw[arrow] (b2) -- node[right] {false} (b4);
\draw[arrow] (b3) -- (b4);
\end{tikzpicture}
\end{center}

Множества:

\[
GEN[B_1]=\{d_1\}, \quad KILL[B_1]=\{d_2\},
\]

\[
GEN[B_3]=\{d_2\}, \quad KILL[B_3]=\{d_1\}.
\]

После вычисления получаем:

\[
IN[B_4] = \{d_1,d_2\}.
\]

Это означает, что в строке \texttt{y = x} значение \texttt{x} может происходить как из определения $d_1$, так и из определения $d_2$.

\subsection{Живые переменные}

\begin{definition}
Переменная $x$ называется \textbf{живой} в точке программы, если существует путь выполнения от этой точки, на котором текущее значение $x$ будет использовано до следующего присваивания $x$.
\end{definition}

Для блока $B$ вводятся:

\begin{itemize}
  \item $USE[B]$ --- переменные, используемые в $B$ до их определения в $B$;
  \item $DEF[B]$ --- переменные, определяемые в $B$;
  \item $IN[B]$ --- переменные, живые на входе в $B$;
  \item $OUT[B]$ --- переменные, живые на выходе из $B$.
\end{itemize}

Уравнения:

\[
OUT[B] = \bigcup_{S \in succ(B)} IN[S],
\]

\[
IN[B] = USE[B] \cup (OUT[B] \setminus DEF[B]).
\]

Это обратный анализ, так как информация распространяется от выходов к входам.

\subsection{Абстрактная интерпретация}

\subsubsection{Идея абстрактной интерпретации}

Абстрактная интерпретация является математической основой многих статических анализаторов. Её основная идея состоит в том, чтобы не вычислять точные состояния программы, а заменять их более грубыми абстракциями.

Например, вместо точного множества значений переменной:

\[
\{ -3, -2, -1, 0, 1, 2, 3 \}
\]

можно хранить только интервал:

\[
[-3,3].
\]

Или вместо точного значения переменной хранить её знак:

\[
\texttt{NEG},\quad \texttt{ZERO},\quad \texttt{POS},\quad \texttt{TOP}.
\]

\begin{definition}
\textbf{Конкретная область} $C$ --- множество точных состояний программы.
\end{definition}

\begin{definition}
\textbf{Абстрактная область} $A$ --- множество приближённых описаний состояний программы.
\end{definition}

\begin{definition}
Функция $\alpha : C \to A$ называется \textbf{абстракцией}, если она сопоставляет конкретному состоянию или множеству состояний его абстрактное описание.
\end{definition}

\begin{definition}
Функция $\gamma : A \to C$ называется \textbf{конкретизацией}, если она сопоставляет абстрактному элементу множество конкретных состояний, которые он описывает.
\end{definition}

Обычно требуется свойство корректности:

\[
c \subseteq \gamma(\alpha(c)).
\]

Это означает, что абстракция не теряет состояний с точки зрения безопасности анализа: она может добавить лишние состояния, но не должна исключить реальные.

\subsubsection{Решётки и неподвижные точки}

Многие статические анализы формулируются как поиск неподвижной точки монотонной функции на решётке.

\begin{definition}
Частично упорядоченное множество $(L,\sqsubseteq)$ называется \textbf{решёткой}, если для любых $a,b \in L$ существуют:
\begin{itemize}
  \item точная верхняя грань $a \sqcup b$;
  \item точная нижняя грань $a \sqcap b$.
\end{itemize}
\end{definition}

\begin{definition}
Полная решётка --- это частично упорядоченное множество, в котором каждое подмножество имеет точную верхнюю и точную нижнюю грань.
\end{definition}

\begin{definition}
Функция $f : L \to L$ называется \textbf{монотонной}, если
\[
x \sqsubseteq y \Rightarrow f(x) \sqsubseteq f(y).
\]
\end{definition}

\begin{definition}
Элемент $x \in L$ называется \textbf{неподвижной точкой} функции $f$, если
\[
f(x)=x.
\]
\end{definition}

\begin{theorem}[Теорема Кнастера--Тарского]
Пусть $(L,\sqsubseteq)$ --- полная решётка, а $f : L \to L$ --- монотонная функция. Тогда множество неподвижных точек функции $f$ непусто и само образует полную решётку. В частности, существуют наименьшая и наибольшая неподвижные точки:
\[
\operatorname{lfp}(f)=\BigSqcap \{x \in L \mid f(x) \sqsubseteq x\},
\]
\[
\operatorname{gfp}(f)=\BigSqcup \{x \in L \mid x \sqsubseteq f(x)\}.
\]
\end{theorem}

\begin{proof}
Докажем существование наименьшей неподвижной точки.

Обозначим множество предподвижных точек:
\[
P = \{x \in L \mid f(x) \sqsubseteq x\}.
\]

Так как $L$ --- полная решётка, множество $P$ имеет точную нижнюю грань:
\[
p = \BigSqcap P.
\]

Покажем, что $p$ является неподвижной точкой.

Поскольку $p$ является нижней гранью $P$, для любого $x \in P$ верно:
\[
p \sqsubseteq x.
\]

Так как $f$ монотонна, из $p \sqsubseteq x$ следует:
\[
f(p) \sqsubseteq f(x).
\]

Но $x \in P$, значит:
\[
f(x) \sqsubseteq x.
\]

Следовательно:
\[
f(p) \sqsubseteq x
\]
для любого $x \in P$. Значит, $f(p)$ также является нижней гранью множества $P$. Так как $p$ --- наибольшая из всех нижних граней $P$, получаем:
\[
f(p) \sqsubseteq p.
\]

Отсюда $p \in P$.

Теперь применим монотонность $f$ к неравенству $f(p) \sqsubseteq p$:
\[
f(f(p)) \sqsubseteq f(p).
\]

Следовательно, $f(p) \in P$. Так как $p$ --- нижняя грань $P$, а $f(p) \in P$, получаем:
\[
p \sqsubseteq f(p).
\]

Итак, одновременно:
\[
f(p) \sqsubseteq p
\]
и
\[
p \sqsubseteq f(p).
\]

Следовательно:
\[
f(p)=p.
\]

Значит, $p$ является неподвижной точкой.

Покажем, что она наименьшая. Пусть $y$ --- любая неподвижная точка $f$. Тогда:
\[
f(y)=y,
\]
следовательно,
\[
f(y) \sqsubseteq y,
\]
то есть $y \in P$. Поскольку $p$ --- нижняя грань $P$, имеем:
\[
p \sqsubseteq y.
\]

Значит, $p$ --- наименьшая неподвижная точка.

Доказательство существования наибольшей неподвижной точки проводится двойственным образом через множество постподвижных точек:
\[
Q=\{x \in L \mid x \sqsubseteq f(x)\}.
\]

Множество всех неподвижных точек образует полную решётку как замкнутая относительно верхних и нижних граней структура, индуцированная полной решёткой $L$.
\end{proof}

\subsubsection{Пример абстрактной интерпретации: интервальный анализ}

Интервальная абстракция описывает возможные значения переменной интервалом:

\[
[l,u],
\]
где $l,u \in \mathbb{Z}\cup\{-\infty,+\infty\}$ и $l \le u$.

Программа:

\begin{lstlisting}[language=C]
int x = 0;
while (x < 10) {
    x = x + 1;
}
\end{lstlisting}

Интуитивно:

\[
x = 0,
\]
после одной итерации:
\[
x = 1,
\]
после двух:
\[
x = 2,
\]
после $n$ итераций:
\[
x = n.
\]

Абстрактный анализ строит последовательность интервалов:

\[
[0,0] \sqsubseteq [0,1] \sqsubseteq [0,2] \sqsubseteq \cdots \sqsubseteq [0,10].
\]

Итоговый инвариант цикла:

\[
0 \le x \le 10.
\]

Если анализатор не учитывает условие выхода достаточно точно, он может получить более грубый инвариант:

\[
0 \le x.
\]

\subsection{Расширение и сужение в абстрактной интерпретации}

Для циклов и рекурсии поиск неподвижной точки может не завершаться, если решётка имеет бесконечные возрастающие цепи. Поэтому используются операции расширения и сужения.

\begin{definition}
Оператор \textbf{расширения} $\nabla$ ускоряет достижение неподвижной точки, заменяя возрастающую последовательность приближений более грубой, но стабилизирующейся последовательностью.
\end{definition}

Для интервалов типичное расширение:

\[
[l_1,u_1] \nabla [l_2,u_2]
=
[l',u'],
\]

где

\[
l' =
\begin{cases}
l_1, & \text{если } l_2 \ge l_1,\\
-\infty, & \text{если } l_2 < l_1,
\end{cases}
\]

\[
u' =
\begin{cases}
u_1, & \text{если } u_2 \le u_1,\\
+\infty, & \text{если } u_2 > u_1.
\end{cases}
\]

Пример:

\[
[0,0] \nabla [0,1] = [0,+\infty].
\]

После расширения может применяться \textbf{сужение}, уточняющее полученный результат:

\[
[0,+\infty] \quad \leadsto \quad [0,10].
\]

\subsection{Символьное выполнение как статико-динамический метод}

\begin{definition}
\textbf{Символьное выполнение} --- это метод анализа, при котором программа выполняется не на конкретных, а на символьных входных значениях, а вдоль каждого пути накапливаются ограничения на эти значения.
\end{definition}

Пример:

\begin{lstlisting}[language=C]
int f(int x) {
    if (x > 0)
        return x + 1;
    else
        return -x;
}
\end{lstlisting}

Пусть входное значение обозначено символом $X$.

Пути:

\[
X > 0 \Rightarrow return\ X+1,
\]

\[
X \le 0 \Rightarrow return\ -X.
\]

Символьное выполнение может использовать SMT-решатель для поиска конкретного входа, проходящего по заданному пути.

Например, для первого пути достаточно выбрать:

\[
X=1.
\]

Для второго:

\[
X=0.
\]

Основная проблема символьного выполнения --- \textbf{взрыв числа путей}. Если в программе $n$ независимых условных операторов, число путей может достигать:

\[
2^n.
\]

\subsection{Формальная верификация и проверка моделей}

\begin{definition}
\textbf{Формальная верификация} --- математическое доказательство того, что программа или модель удовлетворяет заданной спецификации.
\end{definition}

\begin{definition}
\textbf{Проверка моделей} --- автоматизированный метод проверки конечной модели системы относительно формальной спецификации, часто записанной в темпоральной логике.
\end{definition}

Пусть система задана как структура Крипке:

\[
M = (S,S_0,R,L),
\]

где:

\begin{itemize}
  \item $S$ --- конечное множество состояний;
  \item $S_0 \subseteq S$ --- множество начальных состояний;
  \item $R \subseteq S \times S$ --- отношение переходов;
  \item $L : S \to 2^{AP}$ --- разметка состояний атомарными высказываниями.
\end{itemize}

Проверяемое свойство может иметь вид:

\[
AG\ \neg error,
\]

то есть ``на всех путях во всех будущих состояниях не возникает ошибка''.

\section{Динамическая проверка программного кода}

\subsection{Сущность динамической проверки}

\begin{definition}
\textbf{Динамическая проверка программного кода} --- это анализ программы во время её выполнения на конкретных входных данных и в конкретной среде.
\end{definition}

Динамическая проверка включает:

\begin{itemize}
  \item модульное тестирование;
  \item интеграционное тестирование;
  \item системное тестирование;
  \item регрессионное тестирование;
  \item нагрузочное тестирование;
  \item профилирование;
  \item динамический анализ памяти;
  \item динамический анализ гонок;
  \item фаззинг;
  \item выполнение с инструментированием.
\end{itemize}

Главная особенность динамической проверки: она исследует только те поведения программы, которые реально возникли при выполнении.

\subsection{Тестирование}

\begin{definition}
\textbf{Тестовый случай} --- набор входных данных, условий выполнения и ожидаемого результата.
\end{definition}

\begin{definition}
\textbf{Тестовый оракул} --- механизм, позволяющий определить, является ли результат выполнения теста правильным.
\end{definition}

Пример модульного теста:

\begin{lstlisting}[language=Python]
def abs_value(x):
    if x >= 0:
        return x
    return -x

def test_abs_value():
    assert abs_value(5) == 5
    assert abs_value(-5) == 5
    assert abs_value(0) == 0
\end{lstlisting}

\subsection{Уровни тестирования}

\begin{enumerate}
  \item \textbf{Модульное тестирование} --- проверка отдельных функций, классов, модулей.
  \item \textbf{Интеграционное тестирование} --- проверка взаимодействия компонентов.
  \item \textbf{Системное тестирование} --- проверка всей системы целиком.
  \item \textbf{Приёмочное тестирование} --- проверка соответствия ожиданиям заказчика или пользователя.
  \item \textbf{Регрессионное тестирование} --- проверка того, что изменения не нарушили ранее работавшую функциональность.
\end{enumerate}

\subsection{Методы выбора тестов}

\subsubsection{Разбиение на классы эквивалентности}

Идея: множество входов разбивается на классы, внутри которых программа предположительно ведёт себя одинаково.

\begin{example}
Функция принимает возраст пользователя $age$ и разрешает действие, если:

\[
age \ge 18.
\]

Классы эквивалентности:

\[
age < 18,\qquad age \ge 18.
\]

Достаточные представители:

\[
17,\quad 18.
\]
\end{example}

\subsubsection{Анализ граничных значений}

Ошибки часто возникают на границах условий.

Для условия

\[
1 \le x \le 100
\]

полезные тесты:

\[
0,\ 1,\ 2,\ 99,\ 100,\ 101.
\]

\subsubsection{Тестирование таблиц решений}

Если поведение зависит от нескольких булевых условий, строится таблица решений.

\begin{center}
\begin{tabular}{c c c}
\toprule
Условие $A$ & Условие $B$ & Ожидаемое действие \\
\midrule
0 & 0 & $D_1$ \\
0 & 1 & $D_2$ \\
1 & 0 & $D_3$ \\
1 & 1 & $D_4$ \\
\bottomrule
\end{tabular}
\end{center}

\subsection{Критерии покрытия}

\begin{definition}
\textbf{Покрытие кода} --- мера того, какая часть структуры программы была выполнена набором тестов.
\end{definition}

Основные критерии:

\begin{itemize}
  \item покрытие операторов;
  \item покрытие ветвей;
  \item покрытие условий;
  \item покрытие путей;
  \item покрытие функций;
  \item покрытие требований.
\end{itemize}

\subsubsection{Покрытие операторов}

\[
C_{stmt} = \frac{\text{число выполненных операторов}}{\text{общее число операторов}} \cdot 100\%.
\]

\subsubsection{Покрытие ветвей}

\[
C_{branch} = \frac{\text{число выполненных ветвей}}{\text{общее число ветвей}} \cdot 100\%.
\]

Пример:

\begin{lstlisting}[language=C]
int sign(int x) {
    if (x > 0)
        return 1;
    else
        return 0;
}
\end{lstlisting}

Тест $x=5$ покрывает только ветвь \texttt{true}. Для покрытия ветвей нужны, например:

\[
x=5,\qquad x=0.
\]

\subsubsection{Покрытие путей}

Покрытие путей требует выполнения всех возможных путей в графе потока управления. Для программ с циклами число путей может быть бесконечным, поэтому на практике используются ограниченные варианты покрытия путей.

\subsection{Динамический анализ памяти}

Динамические анализаторы памяти обнаруживают ошибки во время выполнения:

\begin{itemize}
  \item выход за границы выделенного блока;
  \item использование после освобождения;
  \item двойное освобождение;
  \item утечки памяти;
  \item использование неинициализированной памяти.
\end{itemize}

Пример ошибки:

\begin{lstlisting}[language=C]
#include <stdlib.h>

int main() {
    int *p = malloc(sizeof(int));
    free(p);
    *p = 10; // use-after-free
    return 0;
}
\end{lstlisting}

Статически такая ошибка может быть обнаружена не всегда, особенно если освобождение и использование указателя находятся в разных функциях и зависят от условий выполнения. Динамический анализатор обнаружит ошибку при выполнении пути, где после \texttt{free(p)} происходит обращение \texttt{*p = 10}.

\subsection{Инструментирование}

\begin{definition}
\textbf{Инструментирование} --- добавление в программу дополнительного кода, собирающего информацию о её выполнении или проверяющего дополнительные условия во время работы.
\end{definition}

Инструментирование может выполняться:

\begin{itemize}
  \item на уровне исходного кода;
  \item на уровне промежуточного представления компилятора;
  \item на уровне байткода;
  \item на уровне бинарного кода;
  \item во время выполнения через динамическую трансляцию.
\end{itemize}

Пример ручного инструментирования:

\begin{lstlisting}[language=C]
int div_safe(int a, int b) {
    if (b == 0) {
        fprintf(stderr, "division by zero\n");
        abort();
    }
    return a / b;
}
\end{lstlisting}

\subsection{Assertions}

\begin{definition}
\textbf{Утверждение} --- логическое условие, которое должно быть истинным в заданной точке программы.
\end{definition}

Пример:

\begin{lstlisting}[language=C]
#include <assert.h>

int sqrt_int(int x) {
    assert(x >= 0);
    // ...
}
\end{lstlisting}

Если \texttt{x < 0}, программа аварийно завершится, что позволит обнаружить нарушение предположения.

Утверждения могут выражать:

\begin{itemize}
  \item предусловия;
  \item постусловия;
  \item инварианты;
  \item внутренние предположения реализации.
\end{itemize}

\subsection{Контрактное программирование}

\begin{definition}
\textbf{Контракт} функции состоит из предусловий, постусловий и инвариантов, задающих обязательства вызывающего кода и самой функции.
\end{definition}

Для функции деления:

\[
divide(a,b)=\frac{a}{b}
\]

предусловие:

\[
b \ne 0.
\]

Постусловие:

\[
result \cdot b = a
\]
для целочисленного деления требует уточнения с учётом остатка:

\[
a = bq + r,\qquad 0 \le |r| < |b|.
\]

\subsection{Фаззинг}

\begin{definition}
\textbf{Фаззинг} --- метод динамической проверки, при котором программа многократно запускается на автоматически генерируемых, мутируемых или специально сконструированных входных данных с целью обнаружения сбоев и нарушений свойств безопасности.
\end{definition}

Основные виды фаззинга:

\begin{itemize}
  \item \textbf{чёрный ящик} --- генерация входов без знания внутренней структуры программы;
  \item \textbf{серый ящик} --- используется ограниченная информация, например покрытие кода;
  \item \textbf{белый ящик} --- используется глубокий анализ программы, например символьное выполнение.
\end{itemize}

\subsubsection{Алгоритм мутационного фаззинга с обратной связью по покрытию}

\textbf{Вход:}
\begin{itemize}
  \item начальный корпус входов $Corpus$;
  \item тестируемая программа $P$;
  \item функция мутации $mutate$;
  \item метрика покрытия $coverage$.
\end{itemize}

\textbf{Выход:}
\begin{itemize}
  \item набор входов, вызывающих сбой;
  \item расширенный корпус интересных входов.
\end{itemize}

\begin{enumerate}
  \item Инструментировать программу $P$ для сбора покрытия.
  \item Запустить $P$ на начальном корпусе входов и вычислить начальное покрытие.
  \item Пока не исчерпан бюджет времени:
  \begin{enumerate}
    \item выбрать вход $s$ из корпуса;
    \item получить новый вход:
    \[
    s' := mutate(s);
    \]
    \item запустить $P$ на $s'$;
    \item если программа аварийно завершилась, сохранить $s'$ как краш-тест;
    \item если $s'$ увеличивает покрытие, добавить $s'$ в корпус;
    \item при необходимости минимизировать интересные входы.
  \end{enumerate}
\end{enumerate}

\subsubsection{Минимальный пример фаззинга}

Пусть программа содержит ошибку:

\begin{lstlisting}[language=C]
#include <stdio.h>
#include <string.h>
#include <stdlib.h>

int main() {
    char buf[8];
    char input[100];

    if (!fgets(input, sizeof(input), stdin))
        return 0;

    if (input[0] == 'A' &&
        input[1] == 'B' &&
        input[2] == 'C') {
        strcpy(buf, input); // возможно переполнение
    }

    return 0;
}
\end{lstlisting}

Обычная случайная генерация может редко получить строку, начинающуюся с \texttt{ABC}. Фаззер с обратной связью по покрытию сохраняет входы, которые проходят всё глубже:

\[
\texttt{A},
\quad
\texttt{AB},
\quad
\texttt{ABC},
\quad
\texttt{ABCxxxxxxxxxxxxxxxx}.
\]

Итоговый вход может вызвать переполнение буфера.

\section{Сравнение статической и динамической проверки}

\begin{center}
\begin{tabular}{p{4cm} p{5cm} p{5cm}}
\toprule
Критерий & Статическая проверка & Динамическая проверка \\
\midrule
Запуск программы & Не требуется & Требуется \\
Входные данные & Обычно не нужны & Обязательны \\
Покрытие путей & Потенциально много путей & Только реально выполненные пути \\
Ложные срабатывания & Возможны часто & Обычно меньше \\
Пропуск ошибок & Возможен из-за приближений & Возможен из-за неполного покрытия тестами \\
Доказательство отсутствия ошибок & Иногда возможно для заданного класса свойств & Обычно невозможно без полного перебора состояний \\
Масштабируемость & Зависит от точности анализа & Зависит от числа тестов и среды выполнения \\
Типичные инструменты & Компиляторные проверки, линтеры, статические анализаторы & Тестовые фреймворки, санитайзеры, фаззеры, профилировщики \\
\bottomrule
\end{tabular}
\end{center}

\section{Звукость, полнота и проблема разрешимости}

\subsection{Звукость и полнота}

\begin{definition}
Анализ называется \textbf{звуковым} относительно свойства $P$, если всякий раз, когда анализ утверждает, что программа удовлетворяет $P$, это утверждение истинно.
\end{definition}

\begin{definition}
Анализ называется \textbf{полным} относительно свойства $P$, если всякий раз, когда программа действительно удовлетворяет $P$, анализ способен это установить.
\end{definition}

В терминах обнаружения ошибок:

\begin{itemize}
  \item звуковый анализ не пропускает ошибок заданного класса, но может выдавать ложные предупреждения;
  \item полный анализ не выдаёт ложных предупреждений, но в общем случае может пропускать ошибки, если не является звуковым.
\end{itemize}

На практике статические анализаторы часто выбирают компромисс между:

\begin{itemize}
  \item точностью;
  \item производительностью;
  \item масштабируемостью;
  \item числом ложных срабатываний;
  \item числом пропущенных ошибок.
\end{itemize}

\subsection{Неразрешимость многих задач анализа программ}

Многие нетривиальные свойства поведения программ невозможно разрешить алгоритмически для всех программ.

\begin{theorem}[Теорема Райса]
Всякое нетривиальное семантическое свойство частично вычислимых функций неразрешимо.
\end{theorem}

\begin{proof}
Докажем через сведение проблемы останова.

Пусть $P$ --- нетривиальное семантическое свойство частично вычислимых функций. Нетривиальность означает, что существуют две частично вычислимые функции $\varphi_a$ и $\varphi_b$ такие, что:
\[
\varphi_a \text{ обладает свойством } P,
\]
а
\[
\varphi_b \text{ не обладает свойством } P.
\]

Предположим противное: существует алгоритм $R$, который по описанию любой программы решает, обладает ли вычисляемая ею функция свойством $P$.

Покажем, что тогда можно решить проблему останова, что невозможно.

Пусть дана произвольная программа $M$ и вход $w$. Построим новую программу $N$ следующим образом.

На входе $x$ программа $N$ выполняет:

\begin{enumerate}
  \item запустить $M$ на входе $w$;
  \item если $M(w)$ останавливается, вычислить $\varphi_a(x)$;
  \item если $M(w)$ не останавливается, то шаг 2 никогда не будет достигнут.
\end{enumerate}

Тогда:

\begin{itemize}
  \item если $M(w)$ останавливается, то $N$ вычисляет функцию $\varphi_a$;
  \item если $M(w)$ не останавливается, то $N$ нигде не останавливается, то есть вычисляет нигде не определённую функцию.
\end{itemize}

Для строгого доказательства выбирают вместо нигде не определённой функции такую функцию $\varphi_b$, не обладающую свойством $P$, и строят программу, которая при останове $M(w)$ ведёт себя как $\varphi_a$, а иначе ведёт себя как $\varphi_b$. Это достигается стандартной техникой параллельного моделирования: программа $N$ сначала имитирует $M(w)$ ограниченное число шагов, а до обнаружения останова ведёт себя как $\varphi_b$.

Тогда:
\[
M(w) \downarrow \Rightarrow \varphi_N \text{ обладает } P,
\]
\[
M(w) \uparrow \Rightarrow \varphi_N \text{ не обладает } P.
\]

Запустив предполагаемый решатель $R$ на программе $N$, можно было бы определить, останавливается ли $M$ на $w$.

Это противоречит неразрешимости проблемы останова. Следовательно, алгоритма $R$ не существует, и свойство $P$ неразрешимо.
\end{proof}

\begin{remark}
Из теоремы Райса следует, что не существует универсального статического анализатора, который для любой программы точно и всегда за конечное время определял бы наличие произвольного нетривиального поведенческого свойства.
\end{remark}

\section{Типичные ошибки, обнаруживаемые статически и динамически}

\subsection{Ошибки, хорошо обнаруживаемые статически}

\begin{itemize}
  \item синтаксические ошибки;
  \item ошибки типов;
  \item использование необъявленных имён;
  \item некоторые случаи недостижимого кода;
  \item потенциальное использование неинициализированной переменной;
  \item нарушения стандартов кодирования;
  \item некоторые SQL-инъекции и XSS по анализу потоков данных;
  \item часть ошибок управления ресурсами;
  \item подозрительные конструкции.
\end{itemize}

Пример:

\begin{lstlisting}[language=C]
int f(int flag) {
    int x;
    if (flag)
        x = 10;
    return x; // возможно использование неинициализированной переменной
}
\end{lstlisting}

Статический анализ может заметить путь:

\[
flag = 0,
\]

на котором переменная \texttt{x} не инициализируется.

\subsection{Ошибки, хорошо обнаруживаемые динамически}

\begin{itemize}
  \item реальные падения программы;
  \item ошибки памяти при конкретном запуске;
  \item гонки данных, проявившиеся при выполнении;
  \item нарушения assertions;
  \item деградация производительности;
  \item утечки ресурсов;
  \item ошибки конфигурации окружения;
  \item несовместимость с конкретной платформой.
\end{itemize}

Пример:

\begin{lstlisting}[language=C]
int a[10];
a[i] = 1;
\end{lstlisting}

Статически значение $i$ может быть неизвестно. Динамически при запуске можно проверить:

\[
0 \le i < 10.
\]

\section{Статическая проверка безопасности}

Статический анализ безопасности часто строится как анализ потоков данных от источников к приёмникам.

\begin{definition}
\textbf{Источник} --- точка программы, из которой поступают потенциально недоверенные данные.
\end{definition}

\begin{definition}
\textbf{Приёмник} --- опасная операция, использование в которой недоверенных данных может привести к уязвимости.
\end{definition}

\begin{definition}
\textbf{Санитайзер} --- операция, очищающая или валидирующая данные так, что они становятся безопасными для определённого приёмника.
\end{definition}

Пример:

\begin{lstlisting}[language=Python]
name = request.GET["name"]      # источник
sql = "SELECT * FROM users WHERE name = '" + name + "'"
cursor.execute(sql)             # приемник
\end{lstlisting}

Анализатор должен отследить поток:

\[
request.GET["name"] \to name \to sql \to cursor.execute(sql).
\]

Если между источником и приёмником нет корректной параметризации запроса или экранирования, возникает предупреждение о SQL-инъекции.

\section{Динамическая проверка безопасности}

Динамическая проверка безопасности включает:

\begin{itemize}
  \item DAST-анализ веб-приложений;
  \item фаззинг сетевых протоколов;
  \item тестирование API;
  \item динамический анализ бинарных файлов;
  \item мониторинг системных вызовов;
  \item песочницы;
  \item runtime-защиту.
\end{itemize}

Пример динамического теста SQL-инъекции:

\[
\texttt{' OR '1'='1}
\]

Если сервер возвращает необычно большое число записей или меняет структуру ответа, это может свидетельствовать о наличии уязвимости.

\section{Комбинированные подходы}

Современные процессы качества используют совместно статические и динамические методы.

\subsection{SAST}

\begin{definition}
\textbf{SAST} --- статический анализ безопасности приложений, выполняемый без запуска программы.
\end{definition}

\subsection{DAST}

\begin{definition}
\textbf{DAST} --- динамический анализ безопасности приложений, выполняемый путём взаимодействия с работающим приложением.
\end{definition}

\subsection{IAST}

\begin{definition}
\textbf{IAST} --- интерактивный анализ безопасности, объединяющий выполнение приложения с внутренним инструментированием и анализом потоков данных.
\end{definition}

\subsection{Гибридный анализ}

Гибридный анализ может сочетать:

\begin{itemize}
  \item статический анализ для поиска потенциальных опасных мест;
  \item символьное выполнение для генерации входов;
  \item фаззинг для практического исследования поведения;
  \item динамическое инструментирование для подтверждения ошибки.
\end{itemize}

Пример конвейера:

\begin{center}
\begin{tikzpicture}[
  node distance=1.2cm,
  block/.style={rectangle, draw, rounded corners, align=center, minimum width=3.3cm, minimum height=0.8cm},
  arrow/.style={-{Latex[length=3mm]}}
]
\node[block] (static) {Статический анализ};
\node[block, right=of static] (targets) {Подозрительные\\участки};
\node[block, right=of targets] (symb) {Символьное\\выполнение};
\node[block, below=of symb] (fuzz) {Фаззинг};
\node[block, left=of fuzz] (confirm) {Динамическое\\подтверждение};

\draw[arrow] (static) -- (targets);
\draw[arrow] (targets) -- (symb);
\draw[arrow] (symb) -- (fuzz);
\draw[arrow] (fuzz) -- (confirm);
\end{tikzpicture}
\end{center}

\section{Практические принципы применения}

\subsection{Принципы статической проверки}

\begin{enumerate}
  \item \textbf{Раннее обнаружение дефектов.} Статический анализ следует запускать как можно раньше: в IDE, при сборке, в CI.
  \item \textbf{Инкрементальность.} Анализ должен быстро обрабатывать изменения, чтобы не мешать разработке.
  \item \textbf{Настройка правил.} Набор проверок должен соответствовать языку, проекту, стандартам и критичности системы.
  \item \textbf{Управление ложными срабатываниями.} Предупреждения нужно классифицировать, подавлять только обоснованно и документировать исключения.
  \item \textbf{Интеграция с код-ревью.} Статический анализ не заменяет ревью, но помогает находить типовые ошибки.
  \item \textbf{Использование спецификаций.} Аннотации, контракты и типы повышают точность анализа.
\end{enumerate}

\subsection{Принципы динамической проверки}

\begin{enumerate}
  \item \textbf{Наличие тестового оракула.} Тест должен определять, является ли результат корректным.
  \item \textbf{Разнообразие входов.} Нужно покрывать нормальные, граничные и ошибочные случаи.
  \item \textbf{Автоматизация.} Тесты должны запускаться регулярно и воспроизводимо.
  \item \textbf{Изоляция.} Модульные тесты должны минимально зависеть от внешней среды.
  \item \textbf{Регрессионность.} Найденная ошибка должна сопровождаться тестом, предотвращающим её повторное появление.
  \item \textbf{Наблюдаемость.} Логи, трассировки, метрики и дампы состояния упрощают диагностику.
  \item \textbf{Инструментирование.} Санитайзеры, профилировщики и сбор покрытия повышают эффективность тестирования.
\end{enumerate}

\section{Минимальный интегральный пример}

Рассмотрим программу:

\begin{lstlisting}[language=C]
#include <stdio.h>

int divide(int a, int b) {
    return a / b;
}

int main() {
    int x, y;
    scanf("%d %d", &x, &y);
    printf("%d\n", divide(x, y));
    return 0;
}
\end{lstlisting}

\subsection{Статическая проверка}

Статический анализатор может построить поток данных:

\[
scanf \to y \to divide(x,y) \to a/b.
\]

В точке деления требуется условие:

\[
b \ne 0.
\]

Но из программы не следует, что $y \ne 0$. Поэтому возможное предупреждение:

\[
\text{``потенциальное деление на ноль''}.
\]

\subsection{Динамическая проверка}

Тест:

\[
x=10,\quad y=2
\]

проходит успешно.

Тест:

\[
x=10,\quad y=0
\]

вызывает ошибку деления на ноль.

\subsection{Исправление}

\begin{lstlisting}[language=C]
#include <stdio.h>
#include <stdlib.h>

int divide(int a, int b) {
    if (b == 0) {
        fprintf(stderr, "division by zero\n");
        exit(1);
    }
    return a / b;
}
\end{lstlisting}

Теперь предусловие проверяется динамически.

\section{Итоговые выводы}

Статическая и динамическая проверка являются двумя основными классами методов обеспечения качества программ.

Статическая проверка:

\begin{itemize}
  \item не требует запуска программы;
  \item способна анализировать множество потенциальных путей;
  \item хорошо подходит для раннего обнаружения ошибок;
  \item может использовать формальные методы;
  \item ограничена неразрешимостью общих свойств программ;
  \item часто выдаёт ложные предупреждения.
\end{itemize}

Динамическая проверка:

\begin{itemize}
  \item анализирует реальные выполнения программы;
  \item хорошо подтверждает наличие ошибки;
  \item зависит от качества тестовых данных;
  \item не доказывает отсутствие ошибок в непокрытых путях;
  \item эффективна в сочетании с инструментированием, санитайзерами и фаззингом.
\end{itemize}

На практике наиболее надёжный подход состоит в совместном использовании:

\[
\text{статический анализ}
+
\text{тестирование}
+
\text{динамическое инструментирование}
+
\text{фаззинг}
+
\text{код-ревью}
+
\text{формальные спецификации}.
\]

\end{document}
